\textbf{Phần 1. Trắc nghiệm lựa chọn}
\setcounterpageref{ex}{0}

\begin{ex} %[5D12N68]
Số gồm: \textit{Bảy triệu, hai trăm linh năm nghìn, sáu trăm ba mươi đơn vị} viết là
\choice
{\True $7\,205\,630$}
{$7\,250\,630$}
{$7\,025\,630$}
{$7\,256\,300$}
\loigiai{
Số đó là $7\,205\,630$.\\
\textbf{Đáp án của câu hỏi là A}
}
\end{ex}

\begin{ex} %[5D12N68]
Trong số thập phân $45{,}678$, chữ số $7$ có giá trị là
\choice
{$7$}
{$0{,}7$}
{$0{,}007$}
{\True $0{,}07$}
\loigiai{
Chữ số 7 ở hàng phần trăm nên có giá trị là $0{,}07$ hay $\dfrac{7}{100}$.\\
\textbf{Đáp án của câu hỏi là D}
}
\end{ex}

\begin{ex} %[5D12H68]
Phân số nào sau đây là phân số tối giản?
\choice
{$\dfrac{6}{9}$}
{\True $\dfrac{8}{15}$}
{$\dfrac{15}{20}$}
{$\dfrac{14}{21}$}
\loigiai{
$\dfrac{8}{15}$ có ƯCLN$(8, 15) = 1$ nên là phân số tối giản.\\
\textbf{Đáp án của câu hỏi là B}
}
\end{ex}

\begin{ex} %[5D12H68]
Phân số $\dfrac{3}{4}$ được viết dưới dạng số thập phân là
\choice
{$0{,}34$}
{$0{,}43$}
{\True $0{,}75$}
{$3{,}4$}
\loigiai{
$\dfrac{3}{4} = 3 : 4 = 0{,}75$.\\
\textbf{Đáp án của câu hỏi là C}
}
\end{ex}

\begin{ex} %[5D12H68]
Hỗn số $3\dfrac{2}{5}$ viết dưới dạng phân số là
\choice
{$\dfrac{6}{5}$}
{\True $\dfrac{17}{5}$}
{$\dfrac{15}{5}$}
{$\dfrac{32}{5}$}
\loigiai{
$3\dfrac{2}{5} = \dfrac{3 \times 5 + 2}{5} = \dfrac{17}{5}$.\\
\textbf{Đáp án của câu hỏi là B}
}
\end{ex}

\begin{ex} %[5D12H68]
Sắp xếp các số sau theo thứ tự từ bé đến lớn: $4{,}05; 4{,}5; 4{,}15; 4{,}005$.
\choice
{$4{,}5 < 4{,}15 < 4{,}05 < 4{,}005$}
{$4{,}05 < 4{,}005 < 4{,}15 < 4{,}5$}
{$4{,}005 < 4{,}15 < 4{,}05 < 4{,}5$}
{\True $4{,}005 < 4{,}05 < 4{,}15 < 4{,}5$}
\loigiai{
Thứ tự: $4{,}005 < 4{,}05 < 4{,}15 < 4{,}5$.\\
\textbf{Đáp án của câu hỏi là D}
}
\end{ex}

\begin{ex} %[5D12V68]
Tìm chữ số $x$ biết $23{,}45 < 23{,}x1 < 23{,}62$:
\choice
{$x = 4$}
{\True $x = 5$ hoặc $x = 6$}
{$x = 5$}
{$x = 6$}
\loigiai{
$23{,}45 < 23{,}51 < 23{,}61 < 23{,}62$ nên $x = 5$ hoặc $x = 6$.\\
\textbf{Đáp án của câu hỏi là B}
}
\end{ex}

\begin{ex} %[5D12V68]
Rút gọn phân số $\dfrac{72}{108}$ về phân số tối giản ta được
\choice
{\True $\dfrac{2}{3}$}
{$\dfrac{3}{4}$}
{$\dfrac{4}{6}$}
{$\dfrac{6}{9}$}
\loigiai{
$\dfrac{72}{108} = \dfrac{72 : 36}{108 : 36} = \dfrac{2}{3}$.\\
\textbf{Đáp án của câu hỏi là A}
}
\end{ex}

\begin{ex} %[5D12V68]
Làm tròn số thập phân $67{,}846$ đến hàng phần mười ta được
\choice
{$67{,}8$}
{$67{,}85$}
{\True $67{,}8$}
{$67{,}9$}
\loigiai{
Chữ số hàng phần trăm là $4 < 5$ nên làm tròn thành $67{,}8$.\\
\textbf{Đáp án của câu hỏi là C}
}
\end{ex}

\begin{ex} %[5D12V68]
Có bao nhiêu số tự nhiên $x$ thỏa mãn $4{,}75 < x < 9{,}2$?
\choice
{\True $5$ số}
{$4$ số}
{$6$ số}
{$7$ số}
\loigiai{
Các số tự nhiên đó là: $5, 6, 7, 8, 9$ (có 5 số).\\
\textbf{Đáp án của câu hỏi là A}
}
\end{ex}

\begin{ex} %[5D12V68]
Tìm phân số tối giản $\dfrac{a}{b}$ lớn hơn $\dfrac{1}{3}$ nhưng bé hơn $\dfrac{1}{2}$ có mẫu số là $12$. Phân số đó là
\choice
{$\dfrac{4}{12}$}
{$\dfrac{6}{12}$}
{\True $\dfrac{5}{12}$}
{$\dfrac{7}{12}$}
\loigiai{
$\dfrac{1}{3} = \dfrac{4}{12} < \dfrac{5}{12} < \dfrac{6}{12} = \dfrac{1}{2}$.\\
\textbf{Đáp án của câu hỏi là C}
}
\end{ex}

\begin{ex} %[5D12C68]
Cho phân số $\dfrac{19}{35}$. Cần bớt ở tử số và thêm vào mẫu số cùng một số tự nhiên $k$ để được phân số mới bằng $\dfrac{1}{2}$. Số tự nhiên $k$ là
\choice
{$2$}
{$3$}
{$4$}
{\True $1$}
\loigiai{
Tổng tử số và mẫu số không đổi: $19 + 35 = 54$.\\
Tử số mới: $54 : (1 + 2) \times 1 = 18$.\\
Số $k = 19 - 18 = 1$.\\
\textbf{Đáp án của câu hỏi là D}
}
\end{ex}

\textbf{Phần 2. Trắc nghiệm đúng sai}
\setcounterpageref{ex}{0}

\begin{ex} %[5D12H68]
Xét tính đúng sai về cấu tạo và so sánh số:
\choiceTF
{\True Số $0$ là số tự nhiên nhỏ nhất.}
{Phân số $\dfrac{12}{15}$ là phân số tối giản.}
{\True $3{,}4 = 3{,}40 = 3{,}400$.}
{\True Trong số $85{,}125$, chữ số $5$ ở phần nguyên có giá trị gấp $1\,000$ lần chữ số $5$ ở phần thập phân.}
\loigiai{
\begin{itemize}
    \item 0 là số tự nhiên nhỏ nhất. Đúng.
    \item $\dfrac{12}{15} = \dfrac{4}{5}$ chưa tối giản. Sai.
    \item Viết thêm chữ số 0 vào bên phải phần thập phân giá trị không đổi. Đúng.
    \item 5 đơn vị gấp 0,005 số lần: $5 : 0{,}005 = 1000$ lần. Đúng.
\end{itemize}
\begin{itemize}
    \item a) \textbf{Đúng}.
    \item b) \textbf{Sai}.
    \item c) \textbf{Đúng}.
    \item d) \textbf{Đúng}.
\end{itemize}
}
\end{ex}

\begin{ex} %[5D12V68]
Cho các số: $2\dfrac{3}{4}; 2{,}7; \dfrac{11}{4}; 2{,}75$. Xét tính đúng sai:
\choiceTF
{\True $2\dfrac{3}{4} = \dfrac{11}{4}$.}
{\True $\dfrac{11}{4} = 2{,}75$.}
{$2{,}7 > 2{,}75$.}
{\True Sắp xếp tăng dần: $2{,}7 < 2{,}75 = 2\dfrac{3}{4} = \dfrac{11}{4}$.}
\loigiai{
\begin{itemize}
    \item $2\dfrac{3}{4} = \dfrac{11}{4}$. Đúng.
    \item $11 : 4 = 2{,}75$. Đúng.
    \item $2{,}7 < 2{,}75$. Khẳng định lớn hơn là Sai.
    \item Thứ tự tăng dần đúng.
\end{itemize}
\begin{itemize}
    \item a) \textbf{Đúng}.
    \item b) \textbf{Đúng}.
    \item c) \textbf{Sai}.
    \item d) \textbf{Đúng}.
\end{itemize}
}
\end{ex}

\textbf{Phần 3. Trả lời ngắn}
\setcounterpageref{ex}{0}

\begin{ex} %[5D12H68]
Viết phân số $\dfrac{3}{8}$ dưới dạng số thập phân.
\shortans{0,375}
\loigiai{
\textbf{Đáp số:} $0{,}375$\\
$3 : 8 = 0{,}375$.
}
\end{ex}

\begin{ex} %[5D12H68]
Rút gọn phân số $\dfrac{48}{60}$ về phân số tối giản $\dfrac{a}{b}$. Tính $a + b$.
\shortans{9}
\loigiai{
\textbf{Đáp số:} $9$\\
$\dfrac{48}{60} = \dfrac{4}{5} \Rightarrow a = 4, b = 5 \Rightarrow a + b = 9$.
}
\end{ex}

\begin{ex} %[5D12V68]
Tìm số tự nhiên $x$ lớn nhất thỏa mãn $x < 15{,}89$.
\shortans{15}
\loigiai{
\textbf{Đáp số:} $15$\\
Số tự nhiên lớn nhất nhỏ hơn 15,89 là $15$.
}
\end{ex}

\begin{ex} %[5D12V68]
Có bao nhiêu chữ số ở phần thập phân của số $12{,}3456$?
\shortans{4}
\loigiai{
\textbf{Đáp số:} $4$\
Có $4$ chữ số ($3, 4, 5, 6$).
}
\end{ex}

\begin{ex} %[5D12V68]
Tìm phân số bằng phân số $\dfrac{2}{5}$ có tổng của tử số và mẫu số bằng $63$. Tìm tử số của phân số đó.
\shortans{18}
\loigiai{
\textbf{Đáp số:} $18$\\
Tử số: $63 : (2 + 5) \times 2 = 18$.
}
\end{ex}

\begin{ex} %[5D12C68]
Viết số thập phân bé nhất có $4$ chữ số khác nhau trong đó có hai chữ số ở phần thập phân.
\shortans{10,23}
\loigiai{
\textbf{Đáp số:} $10{,}23$\\
Số nhỏ nhất có 4 chữ số khác nhau với 2 chữ số thập phân là $10{,}23$.
}
\end{ex}
